%
% 1-elementar.tex -- Elementare Funktionen
%
% (c) 2026 Prof Dr Andreas Müller
%
\section{Elementare Funktionen
\label{buch:intalg:section:elementar}}
\kopfrechts{Elementare Funktionen}%
Der Benutzer eines Computeralgebrasystems erwartet, dass er
als Integranden Funktionen wie
\[
\!\sqrt{ax^2+bx+c},\quad
\sin(x-\pi),\quad
\sqrt{\tan x},\quad
e^{-x^2}
\]
eingeben kann und als Output eine Stammfunktion in einer ähnlichen
Form erhalten kann.
Zum Beispiel liefert Maxima für die dritte Funktion die Stammfunktion
\begin{align*}
\int\!\sqrt{\tan x}\,dx
&=
-\frac1{2^{\frac32}}
\Bigl(
\log \bigl(\tan x+\!\sqrt{2}\sqrt{\tan x}+1\bigr)
-
\log \bigl(\tan x-\!\sqrt{2}\sqrt{\tan x}+1\bigr)
\Bigr)
\\
&\qquad
+\frac1{\!\sqrt{2}}
\Bigl(
\arctan \bigl( \!\sqrt{2}\sqrt{\tan x}+1 \bigr)
+
\arctan \bigl( \!\sqrt{2}\sqrt{\tan x}-1 \bigr)
\Bigr).
\end{align*}
Natürlich ist es recht anspruchsvoll, auf diese Lösung zu kommen.
Allein durch Anwendung heuristischer Integrationsregeln kommt
man wohl nur schwer darauf.
Es braucht also einen algebraischen Prozess, der diese sehr 
komplizierte Stammfunktion schrittweise aufbauen kann.

In der Stammfunktion kommt nicht überraschend die
Funktion $\!\sqrt{\tan x}$
wieder vor, aber auch Polynome in dieser Funktion sowie Logarithmus
und Arkustangens-Funktionen, mit solchen Polynomen als Argument.
\index{arctan}%
Schreibt man $\vartheta(x)=\!\sqrt{\tan x}$, dann wird die Stammfunktion
zu
\begin{align*}
\int\vartheta(x)\,dx
&=
-\frac1{2\!\sqrt{2}}
\Bigl(
\log \bigl(\vartheta(x)^2 +\!\sqrt{2} \vartheta(x) + 1\bigr)
-
\log \bigl(\vartheta(x)^2 -\!\sqrt{2} \vartheta(x) + 1\bigr)
\Bigr)
\\
&\qquad
+\frac1{\!\sqrt{2}}
\Bigl(
\arctan \bigl( \!\sqrt{2}\vartheta(x) + 1 \bigr)
+
\arctan \bigl( \!\sqrt{2}\vartheta(x) - 1 \bigr)
\Bigr).
\end{align*}
Um diese Funktion aufzubauen, wird zusätzlich zur Funktion $\vartheta(x)$
zunächst die Konstante $\!\sqrt{2}$ benötigt.
Damit lassen sich dann die Polynome
\[
\vartheta(x)^2 + \!\sqrt{2}\vartheta(x) + 1,
\quad
\vartheta(x)^2 - \!\sqrt{2}\vartheta(x) + 1,
\quad
\!\sqrt{2}\vartheta(x)+1,
\quad
\!\sqrt{2}\vartheta(x)-1
\]
in $\vartheta(x)$ bilden.
Im nächsten Schritt muss dann der Logarithmus bzw.~die Arkustangensfunktion
angewendet werden.
Schliesslich kann die Stammfunktion als Linearkombination gebildet werden.

Jahrelange Beschäftigung mit algebraischen Umformungen hat die Intuition
gebildet, die in der obenstehenden Analyse der Struktur der Formel
angewendet worden ist.
Um auf diesen Ideen aufbauend einen Integrationsalgorithmus herstellen
zu können, müssen wir die einzelnen Schritt durch wohldefinierte
algebraische Konstruktionen ersetzen.

%
% Funktionenkörper
%
\subsection{Funktionenkörper}
In Kapitel~\ref{chapter:ratint} wurden das Problem studiert,
rationale Funktionen zu integrieren.
Die meisten Operationen bei der Lösung dieses Problems ergaben
wieder rationale Funktionen.
Nur ganz am Schluss traten Terme auf, deren Integrale mit dem
Logarithmus eines Polynoms berechnet werden können.
Dies zeigt, dass wir zum Aufbau von Stammfunktion zumindest die
arithmetischen Grundoperationen auf beliebige Funktionen anwenden können
müssen.
Die natürliche Struktur, in der sich die Integrationstheorie
abspielt, ist daher ein Körper von Funktionen.
Im Kapitel~\ref{chapter:ratint} war dies der Körper der rationalen
Funktionen.
In diesem Abschnitt muss es also darum gehen, diesen Funktionenkörper
massvoll zu erweitern, dass möglichst viele praktisch wichtige
Funktionen darin vorkommen.

%
% Koeffizientenkörper
%
\subsubsection{Koeffizientenkörper}
In Abschnitt~\ref{buch:ratint:section:bernoulli} haben wir als
Funktionenkörper den Körper $\mathbb{C}(x)$ verwendet, dann
aber später eingesehen, dass der Körper $\mathbb{C}$ zu gross ist.
Wenn man mit irrationalen Zahlen algebraisch rechnen muss, muss
man deren algebraischen Eigenschaften kennen.
Die Nullstellen eines Polynoms, die für die Partialbruchzerlegung
benötigt wurden, sind algebraische Element über einem Körper, der
alle Koeffizienten des Nennerpolynoms enthält.
Der Ausgangskörper für die Konstruktion der rationalen Funktionen
entsteht also durch Hinzufügen weiterer Konstanten zu den rationalen
Zahlen.
Die zusätzlichen Konstanten können algebraisch sein, wie die
Wurzel $\!\sqrt{2}$, oder transzendent, wie $\pi$, für die es keine
polynomielle Relation gibt.
Als Zahlenkörper, aus dem die Koeffizienten der Polynome stammen
muss daher immer eine Erweiterung von $\mathbb{Q}$ verwendet werden,
die durch Hinzufügen endlich vieler Elemente $c_1,\dots,c_n$ entstehen.
In der Kette
\[
\underbrace{\mathbb{Q}}_{\displaystyle = K_0}
\subset
\underbrace{\mathbb{Q}(c_1)}_{\displaystyle = K_1}
\subset
\underbrace{\mathbb{Q}(c_1,c_2)}_{\displaystyle = K_2}
\subset
\underbrace{\mathbb{Q}(c_1,c_2,c_3)}_{\displaystyle = K_3}
\subset
\ldots
%\subset
%\underbrace{\mathbb{Q}(c_1,c_2,c_3,\ldots,c_{n-1})}_{\displaystyle = K_{n-1}}
\subset
\underbrace{\mathbb{Q}(c_1,c_2,c_3,\ldots,c_{n-1},c_n)}_{\displaystyle = K_n}
=
K
\]
wird in jedem schritt genau ein neues Element
$c_k$ hinzugefügt, welches transzendent oder algebraisch ist über
dem vorangegangenen Körper.
Falls $c_k$ algebraisch ist, dann gibt es ein Minimalpolynom
$m_k(x) \in K_{k-1}[x]$, für welches $m_k(c_k)=0$ gilt.

Zu den transzendenten Elementen in $K$ gehören insbesondere auch
Parameter in den Ausdrücken.
Wenn ein CAS den quadratischen Ausdruck $ax^2+bx+c$ integrieren soll,
ohne den Koeffizienten $a$, $b$ und $c$ feste Werte zuzuordnen, dann
müssen diese Elemente als transzendent über jedem zugrundeliegenden
Körper betrachtet werden.
Damit sagt man, dass keine weiteren Beziehungen zwischen diesen
Koeffizienten bekannt sind.

%
% Das Funktionsargument
%
\subsubsection{Das Funktionsargument}
In einem Funktionsausdruck spielt ein Symbol die spezielle Rolle
des Funktionsarguments.
Aus algebraischer Sicht verhält sich der Platzhalter für das
Funktionsargument nicht anders wie jeder andere Parameter.
Im Ausdruck $f(x) = ax^2 + bx + c$ werden die Zeichen $a$, $b$ und $c$
wie gewöhnliche Zahlen behandelt, während das Symbol $x$ für das
Funktionsargument steht, für das man nach der Intuition, die man
beim ersten Kontakt mit dem Funktionsbegriff kennenlernt, beliebige
Werte einsetzen kann.
Aus algebraischer Sicht ist $x$ einfach nur ein weiteres transzendentes
Element, um das der Koeffizientenkörper erweitert wird.
Eine besondere Bedeutung erhält $x$ erst, wenn wir, wie das im
nächsten Abschnitt geschehen soll, mit der Ableitung die Konstanten
von der Variablen $x$ unterscheiden können.

\begin{beispiel}
Die rationale Funktion
\[
f(x) 
=
\frac{ax^2 + \!\sqrt{2}x - \sqrt{\pi}}{x^2+\pi}
\]
hat im Zählerpolynom die drei Koeffizienten $a$, $\!\sqrt{2}$ und
$\!\sqrt{\pi}$, im Nenner kommt ausserdem der Koeffizient $\pi$ hinzu.
Die Elemente $a$ und $\pi$ sind transzendent über den rationalen
Zahlen.
Die Elemente
$\!\sqrt{2}$
und
$\!\sqrt{\pi}$
entstehen
durch algebraische Erweiterung von $\mathbb{Q}$ mit dem
Minimalpolynom $m(x)=x^2-2$ bzw.~von $\mathbb{Q}(\pi)$ mit
dem Minimalpolynom $m(x)=x^2-\pi$.
Die Funktion ist daher ein Element im Körper der rationalen
Funktionen
\[
f(x)
\in
\mathbb{Q}(\!\sqrt{2},\pi,\!\sqrt{\pi},a,x).
\]
Man könnte auch versuchen, $\!\sqrt{\pi}$ als transzendentes Element
über $\mathbb{Q}$ zu betrachten.
Dann müsste man aber das Element $\pi$ als das Quadrat von $\!\sqrt{\pi}$
betrachten.
In diesem Fall wäre es zweckmässiger, die Funktion als
\[
f(x)
=
\frac{ax^2+\!\sqrt{2}-c}{x^2+c^2}
\]
mit einem transzendenten Element $c$ zu schreiben, welches für
$\sqrt{\pi}$ steht.
Diese Version macht es aber schwieriger, eine Beziehung zwischen $\pi$ 
und den trigonometrischen Funktionen herzustellen.
\end{beispiel}

Der kleinste Funktionenkörper, der die Konstanten in $K$ und das
Funktionsargument $x$ enthält, ist der Körper $K(x)$ der rationalen
Funktionen mit Koeffizienten in $K$.
Wir bezeichnen ihn auch mit $F=K(x)$.
Im Folgenden verwenden wir die Konvention, Funktionenkörper mit $F$
und Konstantenkörper mit $K$ zu bezeichnen.


%
% Differentialkörper von Funktionen
%
\subsection{Differentialkörper von Funktionen}
In einem Funktionsausdruck für eine differenzierbare Funktion spielt
das Funktionsargument unter allen Symbolen eine besondere Rolle.
Wir gehen davon aus, dass man nach dem Funktionsargument ableiten
kann, weil es ``variabel'' ist.
Nach den Parametern gedenkt man dagegen nicht abzuleiten, da man
diese als ``konstant'' ansehen möchte.

%
% Derivation
%
\subsubsection{Derivation}
Im rein algebraischen Framework eines Funktionenkörpers ist es nicht
sinnvoll, die Definition der Ableitung als Grenzwert
\[
f'(x)
=
\lim_{h\to 0}
\frac{f(x+h)-f(x)}{h}
\]
eines Differenzenquotienten zu verwenden.
Um den Grenzwert zu definieren, wird ein Ordnungsrelation gebraucht,
ein Konzept, welches sich nicht mit rein algebraischen Eigenschaften
erfassen lässt.
Für den Grenzwert wird ausserdem ein Zahlenkörper gebraucht, in dem
alle Grenzwerte existieren, was in $\mathbb{Q}$ nicht gegeben ist.
Der einzige Körper, der vollständig wäre, ist $\mathbb{R}$, der aber
im Computer nicht exakt dargestellt werden kann und daher nicht als
Basis eines CAS definiert werden kann.

Als alternativer Ansatz für die Definition der können die einfachsten
bekannten Rechenregeln der Ableitung verwendet werden.
Die Ableitung ist zum Beispiel eine lineare Abbildung.
Dies allein reicht jedoch nicht, weil die Ableitung auch für 
Potenzen $x^n$ oder für Produkte von Funktionen angewendet werden
soll.
Die Linearität kann darüber nichts aussagen.
Daher braucht es eine zusätzlich Regel für Produkte, die natürlich
die Produktregel für dei Ableitung wiedergeben soll.

\begin{definition}[Derivation]
Eine \emph{Derivation} in einem Funktionenkörper $F$ ist eine 
lineare Abbildung $D\colon F\to F$, die zusätzlich die Produktregel
\[
D(fg) = (Df)g + f(Dg)\qquad\forall f,g\in F
\]
erfüllt.
\index{Derivation}%
\end{definition}

In der Definition steht nichts davon, nach welcher Variablen da
abgeleitet werden soll.
In der Analysis verwendet man die Notation
\[
\frac{\partial}{\partial a},\quad
\frac{\partial}{\partial b},\quad
\frac{\partial}{\partial c},\quad
\frac{\partial}{\partial x},\quad
\frac{\partial}{\partial y}\qquad\text{oder}\qquad
\frac{d}{dx},\quad\text{und}\quad
\frac{d}{dz},
\]
die die Variable, nach der abgeleitet wird, klar erkennbar macht.
Allerdings ist in der Analysis auch die Notation $f'$ für die
abgeleitete Funktion üblich, aus der man die Variable nicht
ablesen kann.

Die spezielle Kennzeichnung des Funktionsargumentes ist aber auch
nicht nötig.
Betrachten wir $x$ als das Funktionsargument, dann ist $f(x)=x$ die
einfachste nichtkonstante Funktion, die wir uns denken können.
Die Ableitung dieser Funktion nach $x$ ist in der klassischen Analysis
ist $f'(x)=1$, wir können daher das Funktionsargument im Funktionenkörper
durch die Relation $Dx=1$ charakterisieren.

\begin{definition}[Differentialkörper]
\index{Differentialkorper@Differentialkörper}%
Ein \emph{Differentialkörper} $K$ ist ein Funktionenkörper mit einer
Derivation.
\end{definition}

Der Hauptsatz der Infinitesimalrechnung besagt, dass das Integral
einer Funktion auch eine Stammfunktion ist.
Die Definition des Integrals ist daher für die Arbeit eines
Computeralgebrasystems nicht relevant, es muss nur eine Stammfunktion
finden können.
Der klassische Begriff der Stammfunktion verwendet die Ableitung.
Für den vorliegenden algebraischen Zusammenhang kann er wie folgt
verallgemeinert werden.

\begin{definition}[Stammfunktion]
Sei $F$ ein Differentialkörper mit Derivation $D$.
Eine Funktion $g\in F$ heisst eine \emph{Stammfunktion} von 
$f\in F$, wenn $Dg=f$.
\index{Stammfunktion}%
\end{definition}

%
% Ableitungsregeln
%
\subsubsection{Ableitungsregeln}
Die Produktregel reicht bereits, um die aus dem Analysisunterricht
bekannten Ableitungsregeln zu reproduzieren.

\begin{enumerate}
%
% Ableitung einer Potenz
%
\item
Ableitung einer Potenz: wir betrachten die Potenz $f^n$ eines
Elementes $f\in K$.
Die Derivationseigenschaft liefert
\begin{align*}
D(f^n)
&=
D(\underbrace{f\cdot f\cdot\ldots\cdot f}_{\displaystyle n})
\\
&=
Df\cdot\ldots f
+
f\cdot (Df)\cdot\ldots \cdot f
+
\ldots
+
f\cdot f \cdot\ldots \cdot Df
=
nf^{n-1}\cdot Df.
\end{align*}
Dies ist die aus der Analysis wohlbekannte Regel für Ableitungen einer
Potenz.
%
% Ableitung von  Konstanten
%
\item
Ableitung der Eins: Die Eins zeichnet sich durch die
Eigenschaft $1^n=1$ aus.
Lässt man die Derivation darauf wirken, erhält man
\[
D(1^n)
=
n1^{n-1}D1
=
D1
\qquad\Rightarrow\qquad
(n-1)D1 = 0.
\]
Da $n-1\ne 1$ ist, folgt $D1=0$, die Ableitung der Eins ist somit $0$.
Da die Derivation $D$ linear ist, gilt für jedes andere Element
$q\in\mathbb{Q}$ sofort $Dq = D(q\cdot 1) = q\cdot D1 = 0$.
%
% Ableitung von 1/f
%
\item
Ableitung von $1/f$: Die reziproke Funktion $1/f$ ist charakterisiert
durch die Eigenschaft $f\cdot (1/f)=1$.
Die Ableitung davon ist
\begin{align*}
&&
D(f\cdot(1/f))
&=
(Df)\cdot (1/f) + f\cdot D(1/f)
=
D1
=
0
\\
&\Rightarrow&
f\cdot D(1/f) &= -\frac{Df}{f}
\\
&\Rightarrow&
D(1/f)
&=
-\frac{Df}{f^2},
\end{align*}
dies ist die Ableitungsregel einer Potenz $f^{-1}$.
%
% Potenzregel für beliebige Exponenten
%
\item
Potenzregel für beliebige Exponenten: Wir wissen bereits,
dass $D(f^n)=nf^{n-1}Df$ für positive $n$.
Für einen negativen Exponenten ist
\[
D(f^{-n})
=
D\biggl(\frac{1}{f^n}\biggr)
=
-
\frac{D(f^n)}{f^{2n}}
=
-\frac{nf^{n-1}Df}{f^{2n}}
=
-
n
f^{-n-1}
Df,
\]
dies ist die erwartete Ableitungsformel für beliebige ganzzahlige 
Exponenten.
%
% Ableitung von Polynomen
%
\item
Ableitungsregel für Polynome:
Sei $p(X) = p(X) = p_0 + p_1X + \dots + p_nX^n \in K[X]$ ein Polynom
mit Koeffizienten im Koeffizientenkörper.
Man beachte, dass $X$ hier nur ein Platzhalter ist, für den eine beliebige
Funktion eingesetzt werden kann. 
$p(X)$ selbst ist nicht im Funktionenkörper enthalten, erst $p(f)$ 
wird ein Element des Funktionenkörpers.
Insbesondere darf man $X$ nicht mit $x$ verwechseln.

Wir möchten die Ableitung von $Dp(f)$ berechnen.
Nach der Potenzregel gilt
\begin{align*}
Dp(f)
&=
D(p_0+p_1f + p_2f^2 + \ldots + p_{n-1}f^{n-1}+p_nf^n)
\\
&=
p_1Df + p_2\cdot 2f Df + \ldots p_{n-1}(n-1)f^{n-2}Df + p_nnf^{n-1}Df
\\
&=
\bigl(
p_1 + 2p_2 f + 3p_3f^2 + \ldots + p_{n-1}(n-1)f^{n-2} + p_nnf^{n-1}
\bigr)
Df.
\intertext{Das Polynom auf der rechten Seite ist auch bekannt als
das abgeleitete Polynom $p'(X)$, so dass man}
&= p'(f)\, Df
\end{align*}
schreiben kann.
Dies ist formal die Kettenregel für Polynome.

Der Spezialfall $f=x$ liefert die Ableitungsformel
\[
Dp(x) = p'(x)\cdot Dx = p'(x),
\]
d.~h.~die für das abgeleitet Polynom eingeführte Notation ist im 
Fall $X=x$ identisch mit der Ableitung mit der Deriviation $D$.
%
% Quotientenregel
%
\item
Quotientenregel:
Der Quotient $f/g$ kann auch als Produkt $f\cdot (1/g)$ geschrieben
werden.
Daher ist die Ableitung davon
\begin{align*}
D\frac{f}{g}
&=
D\biggl(f\cdot\frac1g\biggr)
=
Df\cdot\frac1g + f\cdot D\frac1g
\\
&=
Df\cdot\frac1g + f\cdot \biggl(-\frac{Dg}{g^2}\biggr)
\\
&=
\frac{Df}{g} - \frac{f\cdot Dg}{g^2}
=
\frac{Df\cdot g - f\cdot Dg}{g^2}.
\end{align*}
Dies ist die bekannte Quotientenregel.
\item Ableitung einer rationalen Funktion:
Eine rationale Funktion ist ein Quotient $p[X]/q[X]$ mit
Polynomen $p,q\in K[X]$.
Wir möchten $p(f)/q(f)$ ableiten:
\begin{align*}
D\biggl(\frac{p(f)}{q(f)}\biggr)
&=
\frac{
(Dp(f))\cdot q(f) - p(f)\cdot Dq(f)
}{
q(f)^2
}
\\
&=
\frac{p'(f)\,Df\cdot q(f)-p(f)\cdot q'(f)\, Df}{q(f)^2}
\\
&=
\frac{p'(f)q(f)-p(f) q'(f)}{q(f)^2}
Df.
\end{align*}
Dies ist die Kettenregel für den Polynombruch $p(X)/q(X)$.

Auch in diesem Fall ergibt der Spezialfall $f=x$ die wohlbekannte
Quotientenregel
\[
D\frac{p(x)}{q(x)}
=
\frac{p'(x)q(x)-p(x)q'(x)}{q(x)^2}
\]
für rationale Funktionen.
\end{enumerate}

Diese Beispiele zeigen, dass alle Regeln, die man üblicherweise
zur Berechnung von Ableitungen braucht, allein aus der Linearität
und Eigenschaft der Derivation folgen.
Der Funktionenkörper $F$ mit der Derivation $D$ ist daher ein
vollständiges Modell für die vertrauten algebraischen Eigenschaften
der Ableitung.
Es ist daher auch denkbar, dass das Problem, eine Stammfunktion
zu finden, allein in diesem algebraischen Rahmen gelöst werden
kann.

Der Vollständigkeit halber halten wir hier noch die Definition
des abgeleiteten Polynoms fest.

\begin{definition}[Abgeleitetes Polynom]
\label{buch:intalg:elementar:def:abgeleitetespolynom}
\index{Abgeleites Polynom}%
Ist $p(X)\in K[X]$ ein Polynom mit Koeffizienten in einem
Koeffizientenkörper mit der Darstellung
\[
p(X)
=
p_0 + p_1X + p_2X^2 + \ldots + p_{n-1}X^{n-1} + p_nX^n,
=
\sum_{k=0}^n p_k X^k
\]
dann heisst
\[
p'(X)
=
p_1 + 2p_2 X + 3p_3 X^2 + \ldots + (n-1)p_{n-1}X^{n-2} + np_n X^{n-1}
=
\sum_{k=1}^n kp_k X^{k-1}
\]
das \emph{abgeleitete Polynom}.
\end{definition}

%
% Ableitung eines algebraischen Elementes
%
\subsubsection{Ableitung eines algebraischen Elementes}
Wenn der Funktionenkörper $F$ um ein algebraisches Element $\alpha$
erweitert wird, dann muss auch die Ableitung $D\alpha$ berechnet
werden können, damit $F(\alpha)$ wieder ein Differentialkörper ist.

Sei also $m(X)\in F[X]$ das Minimalpolynom des Elementes $\alpha$,
es gilt also $m(\alpha)=0$ und folglich auch $Dm(\alpha)=0$.
Andererseits ist nach den Berechungsregeln die Ableitung
\begin{align*}
Dm(\alpha)
&=
D(m_0)
+
D(m_1\alpha)
+
D(m_2\alpha^2)
+
D(m_3\alpha^3)
+
\ldots
+
D(m_n\alpha^n)
\\
&=
{\color{darkred}Dm_0}
+
({\color{darkred}Dm_1})\alpha + {\color{blue}m_1}D\alpha
+
({\color{darkred}Dm_2})\alpha^2 + {\color{blue}2m_2 \alpha}D\alpha
%+
%(Dm_3)\alpha^3 + 3m_3 \alpha^2 D\alpha
+
\ldots
+
({\color{darkred}Dm_n})\alpha^n + {\color{blue}nm_n\alpha^{n-1}}D\alpha
\\
&=
{\color{blue}m'(\alpha)}D\alpha
+
({\color{darkred}Dm_0+Dm_1\alpha + \dots + Dm_n\alpha^n}).
\intertext{Da $m(X)$ das Minimalpolynom von $\alpha$ ist, kann $m'(\alpha)$
nicht verschwinden, denn dann wäre $m'$ ein Polynom kleineren Grades,
das auf $\alpha$ verschwindet.
Somit kann man die Gleichung nach $D\alpha$ auflösen und erhält}
D\alpha
&=
-
\frac{{\color{darkred}Dm_0+(Dm_1)\alpha +\ldots+(Dm_n)\alpha^n}}{{\color{blue}m'(\alpha)}}.
\end{align*}
Die Ableitung eines algebraischen Elementes ist also durch das
Minimalpolynom eindeutig bestimmt.

\begin{beispiel}
Die Quadratwurzel $\!\sqrt{f}$ eines Elementes $f\in F$ hat das
Minimalpolynom $m(X) = X^2 -f \in F[X]$ mit der Ableitung $m'(X)=2X$.
Die Ableitung der Koeffizienten ergibt das Polynom $-Df$.
Damit ist die Ableitung der Quadratwurzel
\begin{align*}
D\!\sqrt{f}
&=
\frac{-\color{darkred}Df}{\color{blue}m'(\!\sqrt{f})}
=
-\frac{1}{2\!\sqrt{f}}\cdot Df.
\end{align*}
Dies ist die Kettenregel für die Wurzelfunktion.
\end{beispiel}

\begin{beispiel}
Die $n$-te Wurzel $\sqrt[n]{f}$ eines Elementes $f\in K$ ist algebraisch
mit dem Minimalpolynom $m(X)=X^n -f$ mit Ableitung $m'(X)=nX^{n-1}$.
Dann ist
\begin{align*}
D\!\sqrt[n]{f}
&=
\frac{-\color{darkred}Df}{\color{blue}m'\bigl(\!\sqrt[n]{f}\bigr)}
=
-
\frac{1}{n\bigl(\!\sqrt[n]{f}\bigr)^{n-1}}
\cdot
Df.
\end{align*}
Dies ist die Kettenregel für die $n$-Wurzel.
\end{beispiel}

%
% Konstantenkörper
%
\subsubsection{Konstantenkörper}
Bei der Konstruktion des Funktionskörpers sind wir davon ausgegangen,
dass wir bereits wissen, welche Objekte Konstanten sind.
Die Ableitung $D$ ermöglicht jetzt aber auch, die Konstanten dadurch
zu definieren, dass ihre Ableitung verschwindet.

\begin{definition}[Konstantenkörper]
Sei $F$ ein Differentialkörper.
Ein Element $c\in F$ heisst \emph{konstant} oder eine \emph{Konstante},
\index{konstant}%
\index{Konstante}%
wenn $Dc=0$ ist.
Die Menge
\[
C=\{c\in F\mid Dc=0\}
\]
der Konstanten heisst der \emph{Konstantenkörper} in $F$.
\index{Konstantenkorper@Konstantenkörper}%
\end{definition}

Wir müssen uns noch überlegen, dass $C$ wirklich ein Körper ist.
Dazu müssen wir nachrechnen, dass Summe, Produkt und Quotient
konstanter Elemente aus $C$ wieder in $C$ liegt.
Wir wählen daher $a,b\in C$, d.~h.~$Da=0$ und $Db=0$.
Dann gilt auch
\begin{align*}
D(a\pm b)
&=
Da\pm Db = 0\pm 0 = 0
&&\Rightarrow&
a\pm b&\in C
\\
D(ab)
&=
(Da)\cdot b + a\cdot (Db)
=
0\cdot b + a\cdot 0
=
0
&&\Rightarrow&
ab&\in C
\\
D\frac{a}{b}
&=
\frac{(Da)\cdot b - a\cdot (Db)}{b^2}
=
\frac{0\cdot b - a\cdot 0}{b^2}
=
0.
&&\Rightarrow&
\frac{a}{b}&\in C.
\end{align*}
$C$ ist somit ein Körper.
Der Konstantenkörper kann grösser sein als der ursprüngliche
Koeffizientenkörper.

%
% Exponential- und Logarithmusfunktionen
%
\subsection{Exponential- und Logarithmusfunktionen}
Die Exponential- und Logarithmusfunktionen werden in der klassischen
Analysis mithilfe einer Potenzreihe oder durch die Umkehrfunktion
definiert.
Beide Konzepte stehen im algebraischen Umfeld eines Differentialkörpers
nicht zur Verfügung.
Es muss also eine alternative Methode gefunden werden, diese Funktionen
zu charakterisieren.

Die Potenzreihenmethode zur Lösung von Differentialgleichungen
hat in Kapitel~\ref{chapter:differentialgleichungen} gezeigt,
dass die Potenzreihe von $e^x$ eine Folge der Differentialgleichung
für die Exponentialgleichung ist.
Die Differentialgleichung lässt sich auch in einem Differentialkörper
formulieren, es liegt daher nahe, diese anstelle der Potenzreihe
zur Charakterisierung der Exponentialfunktion zu verwenden.

\begin{definition}[Exponential- und Logarithmusfunktion]
\label{buch:intalg:elementar:def:explog}
Sei $F$ ein Differentialkörper und $f\in F$.
Ein Element $\vartheta\in F$ heisst eine \emph{Exponentialfunktion}
von $f$, wenn $D\vartheta = \vartheta\cdot Df$.
\index{Exponentialfunktion}%
Die Exponentialfunktion wird manchmal auch $\vartheta=e^f$ oder
$\vartheta=\exp f$ geschrieben.
$\vartheta$ heisst ein \emph{Logarithmus} von $f$, wenn
$D\vartheta = Df/\vartheta$ oder $\vartheta\cdot D\vartheta = Df$.
\index{Logarithmus}%
Der Logarithmus wird manchmal auch $\vartheta=\log f$ geschrieben.
\end{definition}

Schon bei der Definition des Derivationsbegriffs haben wir darauf
hingewiesen, dass es für die Derivation keine Kettenregel gibt,
weil in einem Funktionen das Konzept der Funktionszusammensetzung fehlt.
Die Differentialgleichungen der
Definition~\ref{buch:intalg:elementar:def:explog}
sind die Kettenregeln für die Exponential- bzw.~Logarithmusfunktion.
Die Kettenregel ist in diesem Rahmen also weniger eine Eigenschaft
der Deriviation, sondern der Funktionen, mit denen wir arbeiten wollen.

%
% Eindeutigkeit
%
\subsubsection{Eindeutigkeit}
Exponential- und Logarithmusfunktion sind durch die Bedingungen der
Definition~\ref{buch:intalg:elementar:def:explog} nicht eindeutig
festgelegt.
Sind $\vartheta_1$ und $\vartheta_2$ zwei Logarithmen von $f\in F$,
dann ist
\[
D(\vartheta_1-\vartheta_2)
=
\frac{Df}{f}-\frac{Df}{f}
=
0,
\]
die Differenz ist also eine Konstante $\vartheta_1-\vartheta_2\in C$.
Eine Logarithmusfunktion ist also nur bis auf eine Konstante in $C$
bestimmt.
Die Verwendung der Notation $\log f$ suggeriert zwar eine eindeutig
bestimmte Funktion, es kann aber nur eine Menge von Funktionen
$\vartheta_1 + C$ gemeint sein.

Sind andererseits $\varphi_1$ und $\varphi_2$ Exponentialfunktionen
von $f$, dann ist
\[
D
\frac{\varphi_1}{\varphi_2}
=
\frac{D\varphi_1\cdot\varphi_2 - \varphi_1\cdot D\varphi_2}{\varphi_2^2}
=
\frac{
\varphi_1\cdot Df\cdot\varphi_2 - \varphi_1\cdot \varphi_2\cdot Df
}{\varphi_2^2}
=
0,
\]
der Quotient $\varphi_1/\varphi_2$ ist also eine Konstante.
Eine Exponentialfunktion ist also nur bis auf einen konstanten
Faktor in $C$ bestimmt.
Die Notation $e^f$ suggeriert eine eineutig bestimmte Funktion,
gemeint sein kann aber nur eine Menge $C\varphi_1$ von Funktionen in $F$.

%
% Logarithmus als Umkehrfunktion der Exponentialfunktion
%
\subsubsection{Logarithmus als Umkehrfunktion der Exponentialfunktion}
Das Konzept der Zusammensetzung von Funktionen kommt in einem
Differentialkörper nicht vor, so dass die Frage, ob der
Logarithmus die Umkehrfunktion der Exponentialfunktion ist,
nicht sinnvoll gestellt werden kann.
Der Begriff der Umkehrfunktion ist nur sinnvoll, wenn man
\index{Umkehrfunktion}%
Funktionen zusammensetzen kann.

Trotzdem kann man eine Interpretation finden, nach der der Logarithmus
eine Art Umkehrfunktion ist.
Da die Exponentialfunktion und der Logarithmus durch die
Definition über ihre differentiellen Eigenschaften nicht eindeutig 
bestimmt sind, kann auch die Umkehrbeziehung nur im Sinne einer
Differentialbeziehung gelten.

Sei jetzt $f\in F$ und $\vartheta$ eine Exponentialfunktion von $f$.
Nach Definition erfüllt sie $D\vartheta=\vartheta\cdot Df$.
Weiter sei $\varphi$ ein Logarithmus von $\vartheta$, d.~h.~nach
Definition ist $D\varphi=D\vartheta/\vartheta$.
Setzt man die definierende Eigenschaft der Exponentialfunktion
in die Ableitung des Logarithmus ein, erhält man
\[
D\varphi
=
\frac{D\vartheta}{\vartheta}
=
\frac{\vartheta\cdot Df}{\vartheta}
=
Df.
\]
Dies bedeutet, dass $\varphi$ bis auf eine additive Konstante mit
$f$ übereinstimmt.
Mit der konventionellen Schreibweise für Exponentialfunktion und
Logarithmus kann man dies auch als
\[
D\log e^f
=
\frac{De^f}{e^f}
=
\frac{e^f Df}{e^f}
=
Df
\]
schreiben.
Dies bedeutet, dass $\log e^f = f+c$ mit einer Konstante $c\in C$ ist.

Umgekehrt ist
\begin{align*}
D\frac{e^{\log f}}{f}
&=
\frac{(De^{\log f})\cdot f - e^{\log f}\cdot Df}{f^2}
\\
&=
\frac{e^{\log f}D(\log f)\cdot f - e^{\log f}\cdot Df }{f^2}
\\
&=
\frac{e^{\log f} (Df/f)\cdot f - e^{\log f}\cdot Df }{f^2}
=
0,
\end{align*}
$e^{\log f}$ ist also bis auf einen konstanten Faktor die ursprüngliche
Funktion $f$.

%
% Logarithmus und Exponentialgesetze
%
\subsubsection{Logarithmus- und Exponentialgesetze}
Für die algebraische Rechnung sind Regeln nötig, mit denen Ausdrücke
mit Logarithmen und Exponentialfunktionen umgeformt werden können.
Dazu gehören die Logarithmus- und Exponentialgesetze
\index{Logarithmusgesetz}%
\index{Exponentialgesetz}%
\[
\log ab = \log a + \log b
\qquad\text{und}\qquad
e^{a+b}=e^ae^b.
\]
Da weder Logarithmus noch Exponentialfunktion eindeutig festgelegt
sind, kann man so eine einfache Beziehung nicht mehr erwarten.

Sind $f_1,f_2\in F$ Funktionen in $F$ und $\vartheta_1$ und $\vartheta_2$
Exponentialfunktionen von $f_1$ bzw.~$f_2$, dann gilt
\[
D(\vartheta_1\vartheta_2)
=
(D\vartheta_1)\vartheta_2
+
\vartheta_1(D\vartheta_1)
=
\vartheta_1 (Df_1) \vartheta_2
+
\vartheta_1\vartheta_2 (Df_2)
=
\vartheta_1\vartheta_2 D(f_1+f_2).
\]
Das Produkt $\vartheta = \vartheta_1\vartheta_2$ ist daher wegen
\[
D\vartheta = \vartheta D(f_1+f_2)
\]
eine Exponentialfunktion von $f_1+f_2$.
In der konventionellen Notation wird dies als
\[
e^{f_1} e^{f_2}
=
e^{f_1+f_2}
\]
geschrieben.
Das Produkt der Exponentialfunktionen der Faktoren ist also eine
Exponentialfunktion der Summe.
Das Exponentialgesetz ist also in diesem erweiterten Sinn immer noch
gültig.

Seien jetzt umgekehrt $\vartheta_1$ und $\vartheta_2$ Logarithmen von
$f_1$ bzw.~$f_2$.
Dann gilt
\[
D(\vartheta_1+\vartheta_2)
=
\frac{Df_1}{f_1}
+
\frac{Df_2}{f_2}
=
\frac{Df_1\cdot f_2+f_1\cdot Df_2}{f_1f_2}
=
\frac{D(f_1f_2)}{f_1f_2}.
\]
Schreibt man $f=f_1f_2$, dann ist die Summe
$\vartheta=\vartheta_1+\vartheta_2$ wegen
\[
D\vartheta = \frac{Df}{f}
\]
ein Logarithmus des Produktes $f$.
In der konventionellen Notation bedeutet dies, dass
\[
\log (f_1f_2)
=
\log f_1 + \log f_2
\]
ist.
Die Summe der Logarithmen der Faktoren ist also ein Logarithmus
des Produktes.


\input{chapters/070-intalg/14-trigo.tex}
%
% Elementare Funktionen
%
\subsection{Elementare Funktionen und das Integrationsproblem}
Der Prozess der Erweiterung des Differentialkörpers der 
rationalen Funktionen $\mathbb{Q}(x)$ durch algebraische Elemente,
Logarithmen und Exponentialfunktionen reproduziert also die Vorgehensweise
eines Mathematikers, der ein Integralproblem lösen will.
Er beginnt mit der Konstruktion eines Funktionenkörpers, der den
Integranden enthält.
Dazu wird es notwendig sein, zusätzliche Konstanten wie Wurzeln,
$\pi$, $i$ und andere oder Parameter hinzuzufügen.
Ausserdem müssen möglicherweise wiederholt Exponentialfunktionen,
Logarithmen und algebraische Funktionen hinzugefügt werden.
Alle Funktionen, mit denen ein menschlicher Rechner arbeitet, sind
auf diese Art konstruierbar.
Dies rechtfertigt, diesen Funktionen einen eigenen Namen zu geben.

\begin{definition}[Elementare Funktionen]
Sei $D$ eine Derivation in $\mathbb{Q}(x)$ mit $Dx=1$.
Eine Funktion $f(x)$ heisst \emph{elementar}, wenn $f$ in einer
Erweiterung von $\mathbb{Q}(x)$ enthalten ist, die durch hinzufügen
algebraischer Elemente, Logarithmen oder Exponentialfunktionen
entsteht.
\index{elementare Funktion}%
\end{definition}

Mit dem Begriff der elementaren Funktion lässt sich jetzt das
Problem der Integration in geschlossener Form präzis formulieren.
\index{geschlossene Form}%
Ein menschlicher Rechner erwartet, dass er den Differentialkörper,
in dem der Integrand gefunden werden kann, durch weitere Erweiterungen
so vergrössern kann, dass sich auch die Stammfunktion darin finden
lässt.
Die Stammfunktion ist daher auch eine elementare Funktion.
Das Problem der Integration in geschlossener Form ist daher das
Folgende.

\begin{problem}[Integration in geschlossener Form]
\index{Integration in geschlossener Form}%
Gegeben eine elementare Funktion $f$ in einem Differentialkörper
$F$, finde eine elementare Funktion $g$ mit $Dg = f$.
\end{problem}

%
% Eindeutigkeit der Erweiterung des Ableitungsoperators
%
\subsection{Eindeutigkeit der Erweiterung des Ableitungsoperators}
Wir betrachten in diesem Abschnitt einen Differentialkörper $F$ mit
der Deriviation $D$ und stellen die Frage, ob sich der Differentialoperator
$D$ auf eindeutige Weise auf eine Erweiterung $F(\vartheta)$ ausdehnen
lässt.

%
% Transzendente Erweiterungen
%
\subsubsection{Transzendente Erweiterungen}
Sei $\vartheta$ eine transzendente Erweiterung von $F$ und
\[
f(\vartheta) = \frac{a(\vartheta)}{b(\vartheta)}
\]
eine rationale Funktion von $\vartheta$ mit Koeffizienten in $F$, dargestellt
als Bruch von Polynomen $a(\vartheta),b(\vartheta)\in F[\vartheta]$.
Wir müssen zeigen, dass es nur eine mögliche Definition einer Deriviation
$\tilde{D}$ auf solchen Brüchen gibt.
$\tilde{D}$ muss alle Rechenregeln einer Deriviation erfüllen, insbesondere
die Quotientenregel
\begin{align*}
\tilde{D}f(\vartheta)
&=
\frac{
\tilde{D}a(\vartheta)\cdot b(\vartheta)-a(\vartheta)\cdot\tilde{D}b(\vartheta)
}{
b(\vartheta)^2
},
\intertext{die Linearität}
\tilde{D} a(\vartheta)
&=
\tilde{D} \sum_{k=1}^n a_k\vartheta^k
=
\sum_{k=1}^n \tilde{D}(a_k\vartheta^k)
\intertext{und die Produktregel}
&=
\sum_{k=1}^n \tilde{D}a_k\cdot \vartheta^k
+
\sum_{k=1}^n a_k\cdot \tilde{D} \vartheta^k
=
\sum_{k=1}^n Da_k\cdot \vartheta^k
+
\sum_{k=1}^n a_k\cdot
k
\vartheta^{k-1}
\tilde{D}
\vartheta.
\end{align*}
Die Derivation $\tilde{D}$ ist also eindeutig bestimmt, sobald
$\tilde{D}\vartheta$ festgelegt ist.

Für eine logarithmische Erweiterung $\vartheta=\log u$ mit $u\in F$ gilt
$\tilde{D}\vartheta=Du/u$.
Da $Du\in F$ und $u\in F$ ist die Ableitung von $\vartheta$ eindeutig
festgelegt und damit $\tilde{D}$ auf ganz $F(\vartheta)$.
%
Für eine exponentielle Erweiterung $\vartheta=e^u$ mit $u\in F$ gilt
$\tilde{D}\vartheta=\tilde{D}e^u = e^uu =\vartheta u$.
Auch in diesem Fall ist die Ableitung von $\vartheta$ eindeutig
festgelegt.
Es folgt, dass es genau eine Erweiterung der Deriviation $D$ auf
$F$ zu einer Derivation $\tilde{D}$ auf ganz $F(\vartheta)$ gibt.

%
% Algebraische Erweiterungen
%
\subsubsection{Algebraische Erweiterungen}
Sei $\vartheta$ eine algebraische Erweiterung mit dem Minimalpolynom
$m(X) \in F[X]$ mit Koeffizienten in $F$.
Auch in diesem Fall stellt sich die Frage, ob die Ausdehnung $\tilde{D}$
auf $F(\vartheta)$ eindeutig bestimmt ist.
Mit zum Fall einer transzendenten Erweiterung analogen Rechnungen folgt
auch hier, dass die Derivation eindeutig festgelegt ist, sobald
$\tilde{D}\vartheta$ festgelegt wird.
Durch Ableiten der Definitionsgleichung
\begin{align}
\sum_{k=0}^n m_k\vartheta^k
&=
0
\intertext{erhält man die Gleichung}
\sum_{k=0}^n Dm_k\vartheta^k
+
\sum_{k=0}^n m_k k\vartheta^{k-1}\tilde{D}\vartheta
&=
0
\notag
\\
\sum_{k=0}^n Dm_k\vartheta^k
+
\biggl(\sum_{k=0}^n m_k k\vartheta^{k-1}\biggr)\tilde{D}\vartheta
&=
0,
\notag
\intertext{die man nach}
\tilde{D}\vartheta
&=
-\frac{\displaystyle
\sum_{k=0}^n Dm_k \cdot \vartheta^k
}{\displaystyle
\sum_{k=0}^n m_kk\vartheta^{k-1}
}
\label{buch:intalg:elementar:Derweiterung:alg}
\end{align}
auflösen kann.
Dies zeigt, dass $\tilde{D}\vartheta$ eindeutig festgelegt ist.

Die zusätzliche Schwierigkeit in diesem Fall ist, dass es mehrere
Funktionen $\vartheta$ gibt, die die Gleichung $m(\vartheta)=0$ erfüllen.
Zum Beispiel definiert das Minimalpolynom $m(X) = X^2-f$ zwei Wurzeln
$\!\sqrt{f}$ und $-\!\sqrt{f}$.
Wir müssen in diesem Fall zusätzlich zeigen, dass
unsere Konstruktion von $\tilde{D}\vartheta$ nicht von der Wahl der
Lösung von $m(\vartheta)=0$ abhängt.
Die Formel~\eqref{buch:intalg:elementar:Derweiterung:alg} zeigt aber,
dass die Ableitung vollständig durch die Koeffizienten des
Minimalpolynoms $m$ festgelegt wird.
Ersetzt man also $\vartheta$ durch eine andere Nullstelle von $m$,
hat die resultierende Ableitung die gleichen algebraischen Eigenschaften.





